\documentclass[12pt]{article}
%\include{amsfonts}
\usepackage{amssymb,amsmath,amsthm,latexsym,epsfig,euscript,multicol}
\usepackage{enumitem}
\usepackage[utf8x]{inputenc}
\usepackage{booktabs}
\usepackage{babel}[spanish]
\usepackage{multicol}
% \usepackage{enumerate}

% Caracteres especiales
\def\A{\mathbb{A}}
\def\C{\mathbb{C}}
\def \N{\mathbb{N}}
\def \P{\mathbb{P}}
\def \Q{\mathbb{Q}}
\def \R{\mathbb{R}}
\def \Z{\mathbb{Z}}
\def \sen{\textrm{sen}}

\theoremstyle{definition}
\newtheorem{ejer}{Ejercicio}
\newcommand{\bej}{\begin{ejer}}
\newcommand{\fej}{\end{ejer}}

\def\dt{\Delta t}
\def\dx{\Delta x}

\topmargin-2cm \vsize 29.5cm \hsize 21cm
\setlength{\textwidth}{16.75cm}\setlength{\textheight}{23.5cm}
\setlength{\oddsidemargin}{0.0cm}
\setlength{\evensidemargin}{0.0cm}


\begin{document}

%----------------------------------------------------------
%Encabezado:

\centerline{{\small Universidad de Buenos Aires - Facultad de Ciencias Exactas y Naturales - Depto. de Matemática}}

\vskip 0.2cm

\hrule

\vskip 0.2cm

 \centerline{{\bf\Large{\sc Investigación Operativa}}}

 \vskip 0.2cm

 \centerline{\ttfamily Segundo Cuatrimestre 2025}

\vskip 0.2cm

 \hrule

 \bigskip
 \centerline{\bf Práctica 2: Simplex}
 \bigskip

%--------------------------------------------------------- 
% 1
\bej
Resolver gráficamente los siguientes problemas de programación lineal:
\begin{multicols}{2}
\begin{enumerate}[label=\alph*)]
    \item \[
            \begin{array}{rrcl}
            \max & \multicolumn{3}{l}{z = 7x_1+8x_2} \\
            \text{s.a:} & 4x_1+x_2 & \leq & 100 \\
             & x_1+x_2 & \leq & 80 \\
             & x_1 & \leq & 40 \\
             & x& \geq & 0
            \end{array}
            \]
    \item \[
            \begin{array}{rrcl}
            \min & \multicolumn{2}{l}{z = 3x_1+9x_2} \\
            \text{s.a:} & -5x_1+2x_2 & \leq & 30 \\
             & -3x_1+x_2 & \leq & 12 \\
              & x& \geq & 0
            \end{array}
            \]
\end{enumerate}
\end{multicols}
\fej

%--------------------------------------------------------- 
% 2
\bej
Resuelva los siguientes problemas aplicando Simplex:
\begin{multicols}{2}
\begin{enumerate}[label=\alph*)]
    \item \[
            \begin{array}{rrcl}
            \min & \multicolumn{2}{l}{z = 3x_1-x_2-3x_3} \\
            \text{s.a:} & x_1+3x_2-x_3+2x_4 & \leq & 2 \\
             & 2x_1+x_2-4x_3 & \leq & 1 \\
             & -4x_2+2x_3 & \leq & 10 \\
             & x& \geq & 0
            \end{array}
            \]
    \item \[
            \begin{array}{rrcl}
            \min & \multicolumn{2}{l}{z = x_1-2x_3} \\
            \text{s.a:} & 2x_1+x_2 & \leq & 4 \\
             & -x_1+2x_2+x_3 & \leq & 7 \\
             & x_3 & \geq & -3 \\
             & x& \geq & 0
            \end{array}
            \]
    \item     \[
    \begin{array}{rrl}
        \max & \multicolumn{2}{l}{3x_1+2x_2-4x_3}  \\
        \text{s.a:} & x_1+x_2-2x_3 \leq& 4 \\
         & -2x_1 + 3x_3 \geq & -5 \\
         & 2x_1 + x_2 - 3x_3 \leq & 7 \\
         & x_1, x_2\geq & 0 \\
         & x_3  \leq & 0 
    \end{array}
    \]
    \item     \[
    \begin{array}{rrl}
        \max & \multicolumn{2}{l}{5x_1 + 6x_2 + 9x_3 + 8x_4}  \\
        \text{s.a:} & x_1+2x_2+3x_3+x_4 \leq& 5 \\
         & x_1 + x_2 + 2x_3 + 3x_4 \leq & 3 \\
         & x_1, x_2, x_3, x_4 \geq & 0 \\
    \end{array}
    \]
\end{enumerate}
\end{multicols}

\fej

%--------------------------------------------------------- 
% 3
\bej
Resolver los modelos del ejercicio anterior utilizando SCIP. 
\fej

%--------------------------------------------------------- 
% 4
\bej
¿Puede una variable que acaba de dejar la base volver a entrar en el siguiente paso del
algoritmo Simplex?
\fej

%--------------------------------------------------------- 
% 5
%\bej
%Demuestre que el problema de minimizar $c^tx$ sujeto a $Ax = b$ carece de interés porque sobre
%$\{x : Ax = b\}$ no existe el mínimo de $c^tx$ o bien $c^tx$ es constante.
%\fej

%--------------------------------------------------------- 
% 6
\bej
Supongamos que se ha resuelto un problema de programación lineal y se desea incorporar
al planteo una nueva variable no negativa con sus correspondientes datos. ¿Cómo se puede
proceder sin rehacer todos los cálculos?
\fej

%--------------------------------------------------------- 
% 7
\bej
Resuelva los siguientes problemas de Programación Lineal utilizando el método simplex.
\begin{multicols}{2}
\begin{enumerate}[label=\alph*)]
    \item \[
            \begin{array}{rrcl}
            \min & \multicolumn{2}{l}{z = -5x_1-7x_2-12x_3+x_4} \\
            s.a: & 2x_1+3x_2+2x_3+x_4 & \leq & 38 \\
             & 3x_1+2x_2+4x_3-x_4 & \leq & 55 \\
             & x& \geq & 0
            \end{array}
            \]
    \item \[
            \begin{array}{rrcl}
            \max & \multicolumn{2}{l}{z = 5x_1+3x_2+2x_3} \\
            s.a: & 4x_1+5x_2+2x_3+x_4 & \leq & 20 \\
             & 3x_1+4x_2-x_3+x_4 & \leq & 30 \\
             & x& \geq & 0
            \end{array}
            \]
    \item \[
            \begin{array}{rrcl}
            \min & \multicolumn{2}{l}{z = 3x_1-2x_2-4x_3} \\
            s.a: & 4x_1+5x_2-2x_3 & \leq & 22 \\
             & x_1-2x_2+x_3 & \leq & 30 \\
             & x& \geq & 0
            \end{array}
            \]
    \item \[
            \begin{array}{rrl}
            \min & \multicolumn{2}{l}{z = -6x_1-14x_2-13x_3} \\
            s.a: & x_1+4x_2+2x_3  \leq & 48 \\
             & x_1+2x_2+4x_3 \leq & 60 \\
              & x \geq & 0
            \end{array}
            \]
\end{enumerate}
\end{multicols}
\fej

%--------------------------------------------------------- 
\bej
Considere el siguiente problema de programación lineal:
\[
    \begin{array}{rrcl}
    \min & \multicolumn{2}{l}{z = -\dfrac{3}{4}x_1+150x_2-\dfrac{1}{50}x_3+6x_4} \\[1em]
    s.a: & \dfrac{1}{4}x_1-60x_2-\dfrac{1}{25}x_3+9x_4+x_5 & = & 0 \\[1em]
     & \dfrac{1}{2}x_1-90x_2-\dfrac{1}{50}x_3+3x_4+x_6 & = & 0 \\[1em]
     & x_3+x_7 & = & 1 \\
     & x & \geq & 0
    \end{array}
\]
\begin{enumerate}[label=\alph*)]
    \item Verifique que si se usa como criterio elegir la variable con menor índice para entrar a la base cuando hay empate, entonces el algoritmo no termina.
    \item Verifique que $(\frac{1}{25}, 0, 1, 0, \frac{3}{100}, 0, 0)$ es una solución óptima y que su valor en la función objetivo es $z_0 = -\frac{1}{20}$ .
\end{enumerate}
\fej

%--------------------------------------------------------- 
%\bej
%Considere el modelo lineal:
%\[
%    \begin{array}{rrcl}
%        \min & \multicolumn{3}{l}{z = c^t x} \\
%        \text{s.a:} & Ax & = & b \\
%         & e^t x & = & 1 \\
%         & x_i & \geq & 0 \quad \forall\,i\in\{1,\dots,n-1\} \\
%         & \multicolumn{2}{c}{x_n\text{ libre}}
%    \end{array}
%\]
%donde $e=(1,\dots, 1)^t$, $b,c\in\mathbb{R}^n$ y $A$ está definida por:
%\[  A_{ij} =\begin{cases} 1 \quad \text{si $i=j$ o $j=n$} \\
%    0 \quad \text{c.c.}
%    \end{cases}
%\]
%Usar la restricción $e^tx=1$ para eliminar la variable libre. ¿Se podría hacer lo mismo si $x_n$ no fuera libre?
%\fej

%---------------------------------------------------------
%\bej
%Considere el modelo lineal:
%\[
%    \begin{array}{rrl}
%        \min & \multicolumn{2}{l}{z=c^tx}  \\
%        \text{s.a:} & Ax = & b\\
%         & x  \geq & 0
%    \end{array}
%\]
%Formule un modelo equivalente en forma estándar tal que el vector de términos independientes sea cero.
%
%Pista: se puede hacer introduciendo una variable y una restricción adicionales.
%\fej

%---------------------------------------------------------
\bej
 Halle todos los valores del parámetro $\alpha$ tales que las regiones definidas por las siguientes restricciones presenten vértices degenerados:
\begin{multicols}{2}
\begin{enumerate}[label=\alph*)]
    \item \[
        \begin{array}{rcl}
            x_1+x_2 & \leq & 8 \\
            6x_1+x_2 & \leq & 12 \\
            2x_1+x_2 & \leq & \alpha \\
            x & \geq & 0
        \end{array}
        \]
    \item \[
            \begin{array}{rcl}
             \alpha x_1+x_2 & \geq & 1 \\
             2x_1+x_2 & \leq & 6 \\
             -x_1+x_2 & \leq & 5 \\
             x_1+2x_2 & \geq & 6 \\
             x & \geq & 0
            \end{array}
            \]
\end{enumerate}
\end{multicols}
\fej

%---------------------------------------------------------
\bej
Aplique el test de optimalidad para encontrar todos los valores del parámetro $\alpha$ tales que
$x^\ast=(0, 1, 1, 3, 0, 0)^t$ sea una solución óptima del siguiente problema de programación lineal:
\[
    \begin{array}{rrcl}
    \min & \multicolumn{3}{l}{z = -x_1-\alpha^2x_2+2x_3-2\alpha x_4-5x_5+10x_6} \\
    s.a: & -2x_1-x_2+x_4+2x_6 & = & 2 \\
     & 2x_1+x_2+x_3 & = & 2 \\
     & -2x_1-x_3+x_4+2x_5 & = & 2 \\
     & x & \geq & 0
    \end{array}
\]
\fej

%---------------------------------------------------------
\bej
El siguiente diccionario corresponde a alguna iteración del método simplex para un problema de minimización:
\[
\begin{array}{rrrrrrrrr}
    x_3 & = & c & + & 2x_2 & - & dx_4 & - & 2x_6  \\
    x_1 & = & e & - & fx_2 & + & 2x_4 & - & x_6  \\
    x_5 & = & 1 & - &  & + & gx_4 & - & 43x_6  \\\hline
    z & = & -1 & + & ax_2 & + & bx_4 & + & 3x_6  \\
\end{array}
\]
Hallar condiciones sobre $a,b,\dots,g$ para que se cumpla:
\begin{enumerate}[label=\alph*)]
    \item la base actual es óptima.
    \item la base actual es la única base óptima.
    \item la base actual es óptima pero no única.
    \item el problema no está acotado.
    \item que $x_4$ entre en la base y el cambio en la función objetivo sea cero.
\end{enumerate}
\fej

%---------------------------------------------------------
\bej
Usar el método simplex (de dos fases o Método M) para resolver los siguientes problemas de programación lineal:
\begin{multicols}{2}
\begin{enumerate}[label=\alph*)]
    \item \[
            \begin{array}{rrcl}
            \min & \multicolumn{2}{l}{z = x_2-3x_3+2x_5} \\
            \text{s.a:} & x_1+3x_2-x_3+2x_5 & = & 2 \\
             & -2x_2+4x_3+x_4 & = & 12 \\
             & x& \geq & 0
            \end{array}
            \]
    \item \[
            \begin{array}{rrcl}
            \min & \multicolumn{2}{l}{z = x_1+4x_2+x_3} \\
            \text{s.a:} & 2x_1-2x_2+x_3 & = & 4 \\\
             & x_1-x_3 & = & 1 \\
             & x_2, x_3& \geq & 0 \\
             & \multicolumn{2}{r}{x_1 \text{ libre}} &
            \end{array}
            \]
    \item \[
            \begin{array}{rrcl}
            \min & \multicolumn{3}{l}{z = x_1-2x_2-8x_5} \\
            s.a: & -2x_1+x_2+x_3-x_5 & \geq & 4 \\
             & -x_1+2x_2+x_6 & \geq & 7 \\
             & x_3+\dfrac{1}{3}x_6 & \leq & 11 \\
             & 6x_2+x_4 & \leq & 3 \\
             & x & \geq & 0
            \end{array}
            \]
    \item \[
            \begin{array}{rrcl}
            \min & \multicolumn{3}{l}{z = 5x_1+2x_2+3x_3} \\
            s.a: & 5x_1-3x_2+6x_3 & \leq & 30 \\
             & x_1+x_2+x_3 & \geq & 14 \\
             & 3x_1-4x_3 & \geq & -15 \\
             & x_1-2x_3 & = & 0 \\
             & x & \geq & 0
            \end{array}
            \]
    \item \[
            \begin{array}{rrcl}
            \min & \multicolumn{3}{l}{z = -4x_1-2x_2-8x_3} \\
            s.a: & 2x_1-x_2+3x_3 & = & 30 \\
             & x_1+2x_2+4x_3 & = & 40 \\
             & x & \geq & 0
            \end{array}
            \]
    \item \[
            \begin{array}{rrcl}
            \min & \multicolumn{3}{l}{z = 2x_1-2x_2-x_3-2x_4+3x_5} \\
            s.a: & -2x_1+x_2-x_3-x_4 & = & 1 \\
             & x_1-x_2+2x_3+x_4+x_5 & = & 4 \\
             & -x_1+x_2-x_5 & = & 4 \\
             & x & \geq & 0
            \end{array}
            \]

\end{enumerate}
\end{multicols}

\fej

% \bej
% Dado un torneo de $n$ equipos donde juegan todos contra todos entre sí $r$ veces, se desea encontrar la mayor cantidad de puntos que un equipo puede obtener y quedar en una posición menor o igual a $k$ al finalizar el torneo. Para resolver este problema, se plantean dos modelos:

% \textbf{Primer modelo:}
%  \begin{itemize}
%         \item $x_{ij}:$ cantidad de veces que el equipo $i$ le gan\'o al equipo $j$ a lo largo del torneo, $x_{ij}\in\mathbb{Z}_{\geq 0}$
%         \item $P_i:$ puntaje del equipo $i$ al finalizar el torneo, $P_i\in \mathbb{Z}_{\geq 0}$
%         \item $\theta_i:$ vale $1$ si y sólo si $P_i \geq P_1$, $\theta_i\in\{0,1\}$
%     \end{itemize}
%     \[\setlength\arraycolsep{1.5pt}
%   \begin{array}{l r c l l r}
%     \max          & P_1  & & & & \\
%     \mathrm{s.a:} & \sum\limits_{i = 2}^n \theta_i & \geq & k-1 & & \\
%                   & \sum\limits_{j \neq i} \left ( 3x_{ij} +(r-x_{ij}-x_{ji} ) \right ) & = & P_i & & \forall i \in \{1,\dots,n \} \\
%                  & x_{ij} & \leq & r & & \forall (i,j) \in \{1,\dots,n \}^2 \\
%                  & x_{ij} + x_{ji} & \leq & r & & \forall (i,j) \in \{1,\dots,n \}^2 \\
%                  & P_i  & \geq & P_1 - M\cdot (1-\theta_i) & & \forall i \in \{1,\dots,n \} \\
%                  & P_i  & \leq & P_1 + M\cdot \theta_i-1 & & \forall i \in \{1,\dots,n \} \\
%                  & \theta_i & \in & \{0, 1\} & & \forall i \in \{1,\dots,n \} \\
%                  & x_{ij} & \in & \mathbb{Z}_{\geq 0} & & \forall (i,j) \in \{1,\dots,n \}^2 \\
%                  & P_i & \in & \mathbb{Z}_{\geq 0} & & \forall (i,j) \in \{1,\dots,n \}^2
%   \end{array}
%     \]
    
%     \textbf{    Segundo modelo:}
    
%     \begin{itemize}
%         \item $x_{ij}:$ cantidad de veces que el equipo $i$ le gan\'o al equipo $j$ a lo largo del torneo, $x_{ij}\in\mathbb{Z}_{\geq 0}$
%         \item $P_i:$ puntaje del equipo $i$ al finalizar el torneo, $P_i\in \mathbb{Z}_{\geq 0}$
%     \end{itemize}
    
%     \[
%       \begin{array}{l r c l l r}
%     \max          & P_k  & & & & \\
%     \mathrm{s.a:} & P_i & \geq & P_{i+1} & & \forall i \in \{1,\dots,n-1 \} \\
%                   & \sum\limits_{j \neq i} \left ( 3x_{ij} +(r-x_{ij}-x_{ji} ) \right ) & = & P_i & & \forall i \in \{1,\dots,n \} \\
%                   & x_{ij} & \leq & r & & \forall (i,j) \in \{1,\dots,n \}^2 \\
%                  & x_{ij} + x_{ji} & \leq & r & & \forall (i,j) \in \{1,\dots,n \}^2 \\
%                 & x_{ij} & \in & \mathbb{Z}_{\geq 0} & & \forall (i,j) \in \{1,\dots,n \}^2 \\
%                  & P_i & \in & \mathbb{Z}_{\geq 0} & & \forall (i,j) \in \{1,\dots,n \}^2
%   \end{array}
%     \]
% Utilizando SCIP, resolver ambos modelos para cada uno de los siguientes conjuntos de parámetros:
% \begin{enumerate}[label=\alph*)]
%     \item $n=4,\;k=3,\; r=1$
%     \item $n=10,\;k=5,\; r=2$
%     \item $n=20,\;k=5,\; r=2$
%     \item $n=30,\;k=4,\; r=2$
% \end{enumerate}

% ¿Cuál de los dos modelos se resuelve más rápido? ¿Por qué?
% \fej
\end{document}
%--------------------------------------------------------- 